\documentclass[11pt,a4paper]{article}
\usepackage[utf8]{inputenc}
\usepackage[margin=1in]{geometry}
\usepackage{listings}
\usepackage{xcolor}
\usepackage{hyperref}
\usepackage{graphicx}
\usepackage{booktabs}
\usepackage{amsmath}

% Code listing style
\lstdefinestyle{rustcode}{
    language=C,
    basicstyle=\ttfamily\small,
    keywordstyle=\color{blue}\bfseries,
    commentstyle=\color{gray}\itshape,
    stringstyle=\color{red},
    numbers=left,
    numberstyle=\tiny\color{gray},
    stepnumber=1,
    numbersep=5pt,
    backgroundcolor=\color{white},
    showspaces=false,
    showstringspaces=false,
    showtabs=false,
    frame=single,
    tabsize=2,
    captionpos=b,
    breaklines=true,
    breakatwhitespace=false,
    escapeinside={(*@}{@*)},
    morekeywords={let, fn, match, if, else, while, for, return, mut, use, enum, struct, impl}
}

\lstdefinestyle{arcturuscode}{
    basicstyle=\ttfamily\small,
    keywordstyle=\color{blue}\bfseries,
    commentstyle=\color{gray}\itshape,
    numbers=left,
    numberstyle=\tiny\color{gray},
    stepnumber=1,
    numbersep=5pt,
    frame=single,
    breaklines=true,
    morekeywords={label, let, jump, call, return, exit, goto, jumpif, callif, print, input}
}

\title{\textbf{Arcturus Language Toolchain} \\ \large Complete Technical Documentation}
\author{Varan}
\date{\today}

\begin{document}

\maketitle

\begin{abstract}
Arcturus is a block-based, stack-oriented programming language with an explicit control flow model using labels and jumps. This document provides comprehensive technical documentation for the complete Arcturus toolchain, including the Muphrid virtual machine, Spica assembler/disassembler, and Hydrae compiler frontend. The toolchain is implemented in Rust and designed for simplicity, explicit control, and educational purposes.
\end{abstract}

\tableofcontents
\newpage

\section{Introduction}

\subsection{Overview}
The Arcturus language toolchain consists of three primary components:

\begin{itemize}
    \item \textbf{Hydrae}: The frontend compiler that processes Arcturus source code (\texttt{.arc}) through preprocessing, lexical analysis, parsing, and bytecode generation
    \item \textbf{Spica}: The assembler and disassembler that converts between human-readable assembly mnemonics and binary bytecode
    \item \textbf{Muphrid}: The virtual machine and interpreter that executes Arcturus bytecode (\texttt{.avm})
\end{itemize}

\subsection{Design Philosophy}
Arcturus is designed with the following principles:

\begin{enumerate}
    \item \textbf{Explicit Control Flow}: All control flow is explicit through labels and jumps, providing complete visibility into program execution
    \item \textbf{Block-Based Structure}: Code is organized into named blocks with clear entry and exit points
    \item \textbf{Stack-Based Execution}: The VM uses a stack-based architecture for simplicity and ease of understanding
    \item \textbf{No Implicit Behavior}: Everything is visible - no hidden allocations, no implicit conversions, no magic
    \item \textbf{Total Control}: Developers have complete control over execution flow and data management
\end{enumerate}

\subsection{Toolchain Pipeline}

The complete compilation and execution pipeline:

\begin{verbatim}
Source Code (.arc)
    ↓
Hydrae Preprocessor (%scope directives)
    ↓
Combined Source
    ↓
Hydrae Lexer (tokenization)
    ↓
Token Stream
    ↓
Hydrae Parser (AST generation)
    ↓
Abstract Syntax Tree
    ↓
Hydrae Compiler (bytecode generation)
    ↓
Bytecode (.avm)
    ↓
Muphrid VM (execution)
\end{verbatim}

Alternatively, assembly can be written directly and processed through Spica:

\begin{verbatim}
Assembly (.asm)
    ↓
Spica Assembler
    ↓
Bytecode (.avm)
    ↓
Muphrid VM
\end{verbatim}

\newpage
\section{Muphrid: The Virtual Machine}

\subsection{Architecture}

Muphrid is a stack-based virtual machine with the following components:

\begin{itemize}
    \item \textbf{Evaluation Stack}: Primary stack for operands and results
    \item \textbf{Call Stack}: Separate stack for managing function calls and return addresses
    \item \textbf{Variable Storage}: HashMap-based storage for named variables
    \item \textbf{Label Table}: HashMap mapping label names to bytecode addresses
    \item \textbf{Instruction Pointer (IP)}: Current position in the bytecode
\end{itemize}

\subsection{Data Types}

The VM supports four fundamental data types:

\begin{lstlisting}[style=rustcode]
enum StackType {
    Int(i64),      // 64-bit signed integer
    Float(f64),    // 64-bit floating point
    String(String),// UTF-8 string
    Bool(bool),    // Boolean value
}
\end{lstlisting}

All values on the stack and in variables are typed using this enum.

\subsection{Bytecode Format}

\subsubsection{General Structure}

All bytecode files must:
\begin{itemize}
    \item Begin with \texttt{ARC\_START} (\texttt{0xF0})
    \item End with \texttt{ARC\_END} (\texttt{0xF1})
    \item Use little-endian encoding for multi-byte values
    \item Store strings as length-prefixed UTF-8: \texttt{[length\_byte][utf8\_bytes]}
\end{itemize}

\subsubsection{Instruction Set}

The complete Arcturus instruction set organized by category:

\paragraph{Arithmetic Operations}
All arithmetic operations are parameterized - they pop N values from the stack, perform the operation, and push the result.

\begin{table}[h]
\centering
\begin{tabular}{@{}llll@{}}
\toprule
Opcode & Mnemonic & Format & Description \\
\midrule
\texttt{0x01} & ADD & \texttt{[0x01][count]} & Pop N integers, sum, push result \\
\texttt{0x02} & SUB & \texttt{[0x02][count]} & Pop N integers, subtract, push result \\
\texttt{0x03} & MUL & \texttt{[0x03][count]} & Pop N integers, multiply, push result \\
\texttt{0x04} & DIV & \texttt{[0x04][count]} & Pop N integers, divide, push result \\
\texttt{0x05} & MOD & \texttt{[0x05][count]} & Pop N integers, modulo, push result \\
\bottomrule
\end{tabular}
\caption{Arithmetic Instructions}
\end{table}

\paragraph{Stack Operations}

\begin{table}[h]
\centering
\begin{tabular}{@{}llll@{}}
\toprule
Opcode & Mnemonic & Format & Description \\
\midrule
\texttt{0x11} & PUSH\_INT & \texttt{[0x11][8 bytes]} & Push i64 constant \\
\texttt{0x12} & PUSH\_STR & \texttt{[0x12][len][bytes]} & Push string constant \\
\texttt{0x13} & PUSH\_BOOL & \texttt{[0x13][bool]} & Push boolean constant \\
\texttt{0x14} & PUSH\_DEC & \texttt{[0x14][8 bytes]} & Push f64 constant \\
\bottomrule
\end{tabular}
\caption{Push Instructions}
\end{table}

\paragraph{Variable Operations}

\begin{table}[h]
\centering
\begin{tabular}{@{}llll@{}}
\toprule
Opcode & Mnemonic & Format & Description \\
\midrule
\texttt{0x10} & STORE & \texttt{[0x10][len][name]} & Pop value, store in variable \\
\texttt{0x15} & LOAD\_INT & \texttt{[0x15][len][name]} & Load variable, push to stack \\
\texttt{0x16} & LOAD\_STR & \texttt{[0x16][len][name]} & Load variable, push to stack \\
\texttt{0x17} & LOAD\_BOOL & \texttt{[0x17][len][name]} & Load variable, push to stack \\
\texttt{0x18} & LOAD\_DEC & \texttt{[0x18][len][name]} & Load variable, push to stack \\
\bottomrule
\end{tabular}
\caption{Variable Instructions}
\end{table}

\paragraph{Comparison Operations}

\begin{table}[h]
\centering
\begin{tabular}{@{}llll@{}}
\toprule
Opcode & Mnemonic & Format & Description \\
\midrule
\texttt{0xC0} & COMPARE & \texttt{[0xC0][op]} & Pop 2 values, compare, push bool \\
\texttt{0xC3} & EQ & \multicolumn{2}{l}{Equality operator (used with COMPARE)} \\
\texttt{0xC4} & GE & \multicolumn{2}{l}{Greater-than-or-equal operator} \\
\texttt{0xC5} & GT & \multicolumn{2}{l}{Greater-than operator} \\
\texttt{0xC6} & LE & \multicolumn{2}{l}{Less-than-or-equal operator} \\
\texttt{0xC7} & LT & \multicolumn{2}{l}{Less-than operator} \\
\texttt{0xC8} & NE & \multicolumn{2}{l}{Not-equal operator} \\
\bottomrule
\end{tabular}
\caption{Comparison Instructions}
\end{table}

\paragraph{Logical Operations}

\begin{table}[h]
\centering
\begin{tabular}{@{}llll@{}}
\toprule
Opcode & Mnemonic & Format & Description \\
\midrule
\texttt{0xD0} & AND & \texttt{[0xD0]} & Pop 2 bools, AND, push result \\
\texttt{0xD1} & OR & \texttt{[0xD1]} & Pop 2 bools, OR, push result \\
\texttt{0xD2} & NOT & \texttt{[0xD2]} & Pop bool, NOT, push result \\
\texttt{0xD3} & XOR & \texttt{[0xD3]} & Pop 2 bools, XOR, push result \\
\bottomrule
\end{tabular}
\caption{Logical Instructions}
\end{table}

\paragraph{Control Flow}

\begin{table}[h]
\centering
\begin{tabular}{@{}llll@{}}
\toprule
Opcode & Mnemonic & Format & Description \\
\midrule
\texttt{0xD4} & LABEL & \texttt{[0xD4][len][name]} & Label marker (preprocessing only) \\
\texttt{0xE4} & JUMP & \texttt{[0xE4][len][label]} & Unconditional jump to label \\
\texttt{0xE2} & CALL & \texttt{[0xE2][len][label]} & Push return addr, jump to label \\
\texttt{0xE3} & RET & \texttt{[0xE3]} & Pop return addr, jump back \\
\texttt{0xC1} & JUMP\_IF & \texttt{[0xC1][len][label]} & Pop bool, jump if true \\
\texttt{0xC2} & CALL\_IF & \texttt{[0xC2][len][label]} & Pop bool, call if true \\
\bottomrule
\end{tabular}
\caption{Control Flow Instructions}
\end{table}

\paragraph{I/O Operations}

\begin{table}[h]
\centering
\begin{tabular}{@{}llll@{}}
\toprule
Opcode & Mnemonic & Format & Description \\
\midrule
\texttt{0xE0} & PRINT & \texttt{[0xE0]} & Pop value and print to stdout \\
\texttt{0xE1} & INPUT & \texttt{[0xE1]} & Read line from stdin, push string \\
\bottomrule
\end{tabular}
\caption{I/O Instructions}
\end{table}

\paragraph{Structural Markers}

\begin{table}[h]
\centering
\begin{tabular}{@{}llll@{}}
\toprule
Opcode & Mnemonic & Format & Description \\
\midrule
\texttt{0xF0} & ARC\_START & \texttt{[0xF0]} & Program start marker \\
\texttt{0xF1} & ARC\_END & \texttt{[0xF1]} & Program end marker \\
\texttt{0xF2} & ARC\_DELIM & \texttt{[0xF2]} & Block delimiter \\
\texttt{0x00} & NOP & \texttt{[0x00]} & No operation \\
\bottomrule
\end{tabular}
\caption{Structural Markers}
\end{table}

\subsection{Execution Model}

\subsubsection{Preprocessing Pass}

Before execution, the VM performs a preprocessing pass to build the label table:

\begin{enumerate}
    \item Scan the bytecode from start to end
    \item When encountering \texttt{LABEL} (\texttt{0xD4}):
    \begin{itemize}
        \item Read the length byte
        \item Read the label name (length bytes of UTF-8)
        \item Store mapping: \texttt{label\_name → instruction\_pointer}
    \end{itemize}
    \item Continue until reaching \texttt{ARC\_END}
\end{enumerate}

\subsubsection{Execution Loop}

The main execution loop:

\begin{lstlisting}[style=rustcode]
while ip < code.len() {
    match code[ip] {
        opcode => {
            // Decode instruction
            // Execute operation
            // Update instruction pointer
        }
    }
}
\end{lstlisting}

Key behaviors:
\begin{itemize}
    \item Most instructions increment IP by their total size
    \item JUMP/CALL set IP to the target label address
    \item RET pops an address from the call stack and sets IP
    \item ARC\_END immediately terminates execution
\end{itemize}

\subsection{Example Bytecode}

A simple program that prints "Hello, World!":

\begin{verbatim}
F0                          ; ARC_START
D4 04 6D 61 69 6E          ; LABEL "main"
12 0D 48 65 6C 6C 6F ...   ; PUSH_STR "Hello, World!"
E0                          ; PRINT
F1                          ; ARC_END
\end{verbatim}

\newpage
\section{Spica: The Assembler}

\subsection{Overview}

Spica provides bidirectional conversion between assembly mnemonics and binary bytecode. It consists of two modes:

\begin{itemize}
    \item \textbf{Assembly Mode}: Converts \texttt{.asm} text files to \texttt{.avm} bytecode
    \item \textbf{Disassembly Mode}: Converts \texttt{.avm} bytecode to \texttt{.asm} text files
\end{itemize}

\subsection{Assembly Syntax}

Assembly files use the following syntax:

\begin{lstlisting}[style=arcturuscode]
ARC_START

LABEL main
PUSH_INT 42
PUSH_INT 8
ADD 2
PRINT

ARC_END
\end{lstlisting}

Key rules:
\begin{itemize}
    \item One instruction per line
    \item Whitespace separates opcode from operands
    \item Multi-word operands are joined with spaces
    \item Comments are not supported (processed by preprocessor)
\end{itemize}

\subsection{Assembler Implementation}

The assembler works line-by-line:

\begin{enumerate}
    \item Read source file
    \item For each non-empty line:
    \begin{itemize}
        \item Split by whitespace
        \item Match first token to mnemonic
        \item Encode operands based on instruction format
        \item Append bytes to output vector
    \end{itemize}
    \item Write bytecode to output file
\end{enumerate}

\subsection{Encoding Examples}

\paragraph{Simple Instruction}
\begin{verbatim}
PRINT  →  [0xE0]
\end{verbatim}

\paragraph{Parameterized Instruction}
\begin{verbatim}
ADD 3  →  [0x01][0x03]
\end{verbatim}

\paragraph{Integer Constant}
\begin{verbatim}
PUSH_INT 42  →  [0x11][0x2A][0x00][0x00][0x00]
                       [0x00][0x00][0x00][0x00]
\end{verbatim}

\paragraph{String with Length Prefix}
\begin{verbatim}
PUSH_STR hello  →  [0x12][0x05][0x68][0x65]
                          [0x6C][0x6C][0x6F]
\end{verbatim}

\paragraph{Label with Length Prefix}
\begin{verbatim}
LABEL main  →  [0xD4][0x04][0x6D][0x61]
                      [0x69][0x6E]
\end{verbatim}

\subsection{Disassembler}

The disassembler reverses the process:

\begin{enumerate}
    \item Read bytecode file
    \item Initialize instruction pointer to 0
    \item While IP < bytecode length:
    \begin{itemize}
        \item Read opcode byte
        \item Based on opcode, read operands
        \item Decode and format as assembly
        \item Print line
        \item Advance IP
    \end{itemize}
\end{enumerate}

The disassembler produces output that can be re-assembled to produce identical bytecode (verified via file diff).

\newpage
\section{Hydrae: The Frontend Compiler}

\subsection{Overview}

Hydrae is the frontend compiler that transforms Arcturus source code into executable bytecode. It consists of three major phases:

\begin{enumerate}
    \item \textbf{Preprocessing}: Handle \texttt{\%scope} directives and module includes
    \item \textbf{Lexical Analysis}: Convert source text into tokens
    \item \textbf{Parsing}: Build abstract syntax tree from tokens
    \item \textbf{Code Generation}: Emit bytecode from AST
\end{enumerate}

\subsection{Preprocessing}

\subsubsection{Purpose}

The preprocessor handles file inclusion before lexical analysis, allowing modular code organization through \texttt{\%scope} directives.

\subsubsection{Scope Directive Syntax}

\begin{verbatim}
%scope lib/math.arc

{main
    call sqrt
end}
\end{verbatim}

\subsubsection{Preprocessing Algorithm}

\begin{enumerate}
    \item Read source file
    \item For each line:
    \begin{itemize}
        \item If line starts with \texttt{\%scope}:
        \begin{itemize}
            \item Extract file path
            \item Read the referenced file
            \item Append content to output
        \end{itemize}
        \item Otherwise, keep line as-is
    \end{itemize}
    \item Append auto-generated exit block:
    \begin{verbatim}
{ label HYD_EXITBLOCK; exit; label HYD_EXITBLOCK_END; }
    \end{verbatim}
    \item Return combined source text
\end{enumerate}

\subsubsection{Design Rationale}

\begin{itemize}
    \item \textbf{Text-level operation}: Preprocessing works on raw text before tokenization, making it simple and efficient
    \item \textbf{No label conflicts}: Since \texttt{\%scope} lines appear outside blocks, they are ignored by the lexer
    \item \textbf{Single output file}: All code is merged into one bytecode file for simplicity
    \item \textbf{Automatic exit}: The HYD\_EXITBLOCK ensures programs terminate cleanly
\end{itemize}

\subsection{Lexical Analysis}

\subsubsection{Token Types}

The lexer produces seven token types:

\begin{lstlisting}[style=rustcode]
enum Token {
    Punct(String),       // Punctuation: (, ), ;, {, }, etc.
    Keyword(String),     // Keywords: let, jump, call, etc.
    Identifier(String),  // Variable/label names
    Integer(i64),        // Integer literals
    Float(f64),          // Floating point literals
    String(String),      // String literals
    Bool(bool),          // Boolean literals: true, false
}
\end{lstlisting}

\subsubsection{Lexer States}

The lexer uses a state machine with six states:

\begin{lstlisting}[style=rustcode]
enum LexerState {
    Idle,       // Inside block, neutral
    Inactive,   // Outside block or in comment
    Punct,      // Building multi-char operator
    Word,       // Building identifier/keyword
    String,     // Inside string literal
    Number,     // Building numeric literal
}
\end{lstlisting}

\subsubsection{State Transitions}

\paragraph{From Idle State:}
\begin{itemize}
    \item \texttt{\}} → Inactive (exit block)
    \item \texttt{\#} → Inactive (enter comment)
    \item \texttt{+, -, *, /, \%, =, !, <, >, \&, |, \^} → Punct
    \item \texttt{;, ,, (, ), \{, [, ], :} → Push single-char punct, stay Idle
    \item \texttt{"} → String
    \item Digit → Number
    \item Letter → Word
    \item Whitespace → Ignore, stay Idle
\end{itemize}

\paragraph{From Punct State:}
Accumulates consecutive operator characters to form multi-character operators like \texttt{==}, \texttt{\&\&}, \texttt{||}, \texttt{\textasciicircum\textasciicircum}.

Exits when encountering:
\begin{itemize}
    \item Whitespace → Push operator, return to Idle
    \item Letter/digit → Push operator, enter Word/Number
    \item Single-char punct → Push operator, push single char
\end{itemize}

\paragraph{From Word State:}
Accumulates alphanumeric characters and underscores.

Exits when encountering non-word character:
\begin{itemize}
    \item Check if word is keyword → Push Keyword
    \item Check if word is \texttt{true}/\texttt{false} → Push Bool
    \item Otherwise → Push Identifier
\end{itemize}

\paragraph{From String State:}
Accumulates all characters until closing \texttt{"}.

\paragraph{From Number State:}
Accumulates digits and optional decimal point.

Exits on non-numeric character:
\begin{itemize}
    \item Contains \texttt{.} → Push Float
    \item Otherwise → Push Integer
\end{itemize}

\paragraph{From Inactive State:}
Ignores all content until:
\begin{itemize}
    \item \texttt{\{} → Enter block, return to Idle
    \item \texttt{\#} → End comment, return to Idle
\end{itemize}

\subsubsection{Example Tokenization}

Input:
\begin{verbatim}
{main
    let x = 42;
    print x;
end}
\end{verbatim}

Token stream:
\begin{verbatim}
[Punct("{"), Keyword("label"), Identifier("main"),
 Keyword("let"), Identifier("x"), Punct("="), Integer(42), Punct(";"),
 Keyword("print"), Identifier("x"), Punct(";"),
 Punct("}")]
\end{verbatim}

\subsection{Parsing}

\subsubsection{AST Node Types}

The parser builds an abstract syntax tree using these node types:

\begin{lstlisting}[style=rustcode]
enum Node {
    Program(Vec<Node>),              // Root node
    Label(String),                   // Label definition
    Exit,                            // Exit instruction
    Goto(String),                    // Unconditional jump
    Jump(String),                    // Jump alias
    Return(Option<Box<Node>>),       // Return from function
    Call(String, Vec<Node>),         // Function call
    JumpIf(Box<Node>, String),       // Conditional jump
    CallIf(Box<Node>, String, Vec<Node>), // Conditional call
    Let(String, Box<Node>),          // Variable assignment
    // Additional nodes for expressions...
}
\end{lstlisting}

\subsubsection{Parsing Strategy}

The parser uses recursive descent with the following structure:

\begin{enumerate}
    \item \texttt{parse\_program()}: Entry point, parses entire token stream
    \item \texttt{parse\_block()}: Parses \texttt{\{...}\} block structure
    \item \texttt{parse\_statement()}: Dispatches based on keyword
    \item \texttt{parse\_expression()}: Handles operator precedence
\end{enumerate}

\subsubsection{Statement Parsing}

Each keyword corresponds to a parsing function:

\begin{lstlisting}[style=rustcode]
match current_token {
    Keyword("label") => parse_label(),
    Keyword("let") => parse_assignment(),
    Keyword("jump") | Keyword("goto") => parse_jump(),
    Keyword("call") => parse_call(),
    Keyword("jumpif") => parse_conditional_jump(),
    Keyword("callif") => parse_conditional_call(),
    Keyword("return") => parse_return(),
    Keyword("exit") => Node::Exit,
    // ...
}
\end{lstlisting}

\subsubsection{Expression Parsing}

Expressions use precedence climbing:

\begin{enumerate}
    \item \texttt{parse\_expression()}: Entry point, calls \texttt{parse\_addition()}
    \item \texttt{parse\_addition()}: Handles \texttt{+} and \texttt{-}
    \item \texttt{parse\_multiplication()}: Handles \texttt{*}, \texttt{/}, \texttt{\%}
    \item \texttt{parse\_primary()}: Handles literals, variables, parentheses
\end{enumerate}

Precedence (low to high):
\begin{enumerate}
    \item Logical: \texttt{||}, \texttt{\&\&}, \texttt{\textasciicircum\textasciicircum}
    \item Comparison: \texttt{==}, \texttt{!=}, \texttt{>}, \texttt{<}, \texttt{>=}, \texttt{<=}
    \item Addition: \texttt{+}, \texttt{-}
    \item Multiplication: \texttt{*}, \texttt{/}, \texttt{\%}
    \item Unary: \texttt{-}, \texttt{!}
    \item Primary: literals, variables, \texttt{( )}
\end{enumerate}

\subsection{Code Generation}

\subsubsection{AST to Bytecode}

The compiler walks the AST and emits bytecode:

\begin{lstlisting}[style=rustcode]
fn compile_node(node: &Node) -> Vec<u8> {
    match node {
        Node::Label(name) => {
            // Emit LABEL instruction
            let mut bytes = vec![0xD4];
            bytes.push(name.len() as u8);
            bytes.extend_from_slice(name.as_bytes());
            bytes
        }
        Node::Let(var, expr) => {
            // Compile expression, then STORE
            let mut bytes = compile_node(expr);
            bytes.push(0x10); // STORE
            bytes.push(var.len() as u8);
            bytes.extend_from_slice(var.as_bytes());
            bytes
        }
        // ... other node types
    }
}
\end{lstlisting}

\subsubsection{Expression Compilation}

Expressions compile to stack-based bytecode:

\paragraph{Example: \texttt{x + y * 2}}

AST:
\begin{verbatim}
BinaryOp(+,
    Variable("x"),
    BinaryOp(*,
        Variable("y"),
        Literal(2)
    )
)
\end{verbatim}

Generated bytecode:
\begin{verbatim}
LOAD_INT x      ; Load x onto stack
LOAD_INT y      ; Load y onto stack
PUSH_INT 2      ; Push 2 onto stack
MUL 2           ; Multiply y * 2
ADD 2           ; Add x + (y * 2)
\end{verbatim}

The compiler performs a post-order traversal (left, right, root) to generate correct stack-based code.

\newpage
\section{Language Reference}

\subsection{Arcturus Source Syntax}

\subsubsection{Block Structure}

All code must be inside blocks:

\begin{lstlisting}[style=arcturuscode]
{block_name
    # Statements go here #
end}
\end{lstlisting}

Content outside blocks is ignored by the lexer.

\subsubsection{Comments}

Comments use hash delimiters:

\begin{lstlisting}[style=arcturuscode]
# This is a comment #

{main
    # Comments can appear anywhere #
    let x = 5;  # Including inline #
end}
\end{lstlisting}

\subsubsection{Variable Declaration}

Variables are declared with \texttt{let}:

\begin{lstlisting}[style=arcturuscode]
let x = 42;
let name = "Alice";
let flag = true;
let pi = 3.14159;
\end{lstlisting}

Variables are dynamically typed and globally scoped.

\subsubsection{Labels and Jumps}

Define labels with \texttt{label}:

\begin{lstlisting}[style=arcturuscode]
{main
    label start;
    print "Hello";
    jump start;  # Infinite loop
end}
\end{lstlisting}

\subsubsection{Conditional Jumps}

Use \texttt{jumpif} for conditional control flow:

\begin{lstlisting}[style=arcturuscode]
let x = 10;
let y = 5;

jumpif x > y bigger_label;
print "x is not bigger";
jump end_label;

label bigger_label;
print "x is bigger";

label end_label;
\end{lstlisting}

\subsubsection{Function Calls}

Use \texttt{call} and \texttt{return}:

\begin{lstlisting}[style=arcturuscode]
{main
    call helper;
    exit;
end}

{helper
    print "Helper called";
    return;
end}
\end{lstlisting}

\subsubsection{I/O Operations}

Print to stdout:
\begin{lstlisting}[style=arcturuscode]
print "Hello, World!";
print x;
\end{lstlisting}

Read from stdin:
\begin{lstlisting}[style=arcturuscode]
input;
let name = input;
\end{lstlisting}

\subsection{Complete Example Program}

\begin{lstlisting}[style=arcturuscode]
# Factorial calculator #

{main
    label main_start;
    
    print "Enter a number: ";
    let n = input;
    
    let result = 1;
    let counter = 1;
    
    label loop;
    jumpif counter > n end_loop;
    
    let result = result * counter;
    let counter = counter + 1;
    jump loop;
    
    label end_loop;
    print "Factorial: ";
    print result;
    
    exit;
end}
\end{lstlisting}

\newpage
\section{Implementation Details}

\subsection{Stack-Based Arithmetic}

Arithmetic operations are parameterized for flexibility:

\begin{lstlisting}[style=rustcode]
0x01 => { // ADD
    let count = code[ip + 1];
    let mut sum = 0i64;
    for _ in 0..count {
        match stack.pop().unwrap() {
            Int(n) => sum += n,
            _ => panic!("ADD expects integers")
        }
    }
    stack.push(Int(sum));
    ip += 2;
}
\end{lstlisting}

This allows efficient multi-operand operations:
\begin{verbatim}
PUSH_INT 1
PUSH_INT 2
PUSH_INT 3
PUSH_INT 4
ADD 4        # Sums all 4 values
\end{verbatim}

\subsection{Label Resolution}

Labels are resolved in a preprocessing pass:

\begin{lstlisting}[style=rustcode]
for _line in &code {
    if np >= code.len() { break; }
    match code[np] {
        0xD4 => {  // LABEL
            let label_name_bytes = &code[np + 2..np + 2 + code[np + 1] as usize];
            let label_name = String::from_utf8(label_name_bytes.to_vec()).unwrap();
            labels.insert(label_name, np);
            np += 2 + code[np + 1] as usize;
        }
        _ => { np += 1 }
    }
}
\end{lstlisting}

This allows O(1) label lookup during execution.

\subsection{Call Stack Management}

Function calls use a separate call stack:

\begin{lstlisting}[style=rustcode]
0xE2 => { // CALL
    let len = code[ip + 1] as usize;
    let label_bytes = &code[ip + 2..ip + 2 + len];
    let label = String::from_utf8(label_bytes.to_vec()).unwrap();
    
    // Push return address
    callstack.push((ip + 2 + len) as i64);
    
    // Jump to function
    ip = *labels.get(&label).unwrap();
}

0xE3 => { // RET
    // Pop and jump to return address
    ip = callstack.pop().unwrap() as usize;
}
\end{lstlisting}

This keeps the call stack separate from the evaluation stack, preventing stack corruption.

\subsection{Variable Storage}

Variables use a HashMap with StackType values:

\begin{lstlisting}[style=rustcode]
let mut data: HashMap<String, StackType> = HashMap::new();

// STORE
0x10 => {
    let var_name = /* decode from bytecode */;
    let value = stack.pop().unwrap();
    match value {
        Int(n) => data.insert(var_name, Int(n)),
        Float(f) => data.insert(var_name, Float(f)),
        String(s) => data.insert(var_name, String(s)),
        Bool(b) => data.insert(var_name, Bool(b)),
    };
}

// LOAD
0x15 => {
    let var_name = /* decode from bytecode */;
    stack.push(data.get(&var_name).unwrap().clone());
}
\end{lstlisting}

All variables are global - there is no scoping mechanism.

\newpage
\section{Usage Guide}

\subsection{Command-Line Interface}

\subsubsection{Hydrae (Compiler)}

Compile Arcturus source to bytecode:
\begin{verbatim}
hydrae program.arc
\end{verbatim}

Output: \texttt{program.arc.avm}

Debug mode (print preprocessing and tokens):
\begin{verbatim}
hydrae program.arc --debug
\end{verbatim}

Lex-only mode (output tokens to file):
\begin{verbatim}
hydrae program.arc --lex
\end{verbatim}

\subsubsection{Spica (Assembler)}

Assemble to bytecode:
\begin{verbatim}
spica assemble program.asm output.avm
\end{verbatim}

Disassemble bytecode:
\begin{verbatim}
spica disassemble program.avm output.asm
\end{verbatim}

\subsubsection{Muphrid (VM)}

Execute bytecode:
\begin{verbatim}
muphrid program.avm
\end{verbatim}

Debug mode (print VM state at each instruction):
\begin{verbatim}
muphrid program.avm -debug
\end{verbatim}

\subsection{Development Workflow}

\subsubsection{Writing High-Level Code}

\begin{enumerate}
    \item Write Arcturus source (\texttt{.arc} file)
    \item Compile with Hydrae
    \item Run with Muphrid
\end{enumerate}

Example:
\begin{verbatim}
$ cat hello.arc
{main
    print "Hello, World!";
    exit;
end}

$ hydrae hello.arc
$ muphrid hello.arc.avm
Hello, World!
\end{verbatim}

\subsubsection{Writing Assembly}

\begin{enumerate}
    \item Write assembly (\texttt{.asm} file)
    \item Assemble with Spica
    \item Run with Muphrid
\end{enumerate}

Example:
\begin{verbatim}
$ cat hello.asm
ARC_START
LABEL main
PUSH_STR Hello, World!
PRINT
ARC_END

$ spica assemble hello.asm hello.avm
$ muphrid hello.avm
Hello, World!
\end{verbatim}

\subsubsection{Debugging}

Use Spica to disassemble compiled code:
\begin{verbatim}
$ hydrae program.arc
$ spica disassemble program.arc.avm debug.asm
$ cat debug.asm
# View generated assembly
\end{verbatim}

Use Muphrid debug mode to trace execution:
\begin{verbatim}
$ muphrid program.avm -debug
ip: 0, opcode: 0xf0
Stack: []
Data: {}
Callstack: []
-----------------------------
# Full execution trace
\end{verbatim}

\newpage
\section{Appendix}

\subsection{Opcode Quick Reference}

\begin{table}[h]
\centering
\small
\begin{tabular}{@{}lll@{}}
\toprule
Category & Opcode & Mnemonic \\
\midrule
\multirow{5}{*}{Arithmetic} & 0x01 & ADD \\
& 0x02 & SUB \\
& 0x03 & MUL \\
& 0x04 & DIV \\
& 0x05 & MOD \\
\midrule
\multirow{5}{*}{Stack} & 0x11 & PUSH\_INT \\
& 0x12 & PUSH\_STR \\
& 0x13 & PUSH\_BOOL \\
& 0x14 & PUSH\_DEC \\
\midrule
\multirow{5}{*}{Variables} & 0x10 & STORE \\
& 0x15 & LOAD\_INT \\
& 0x16 & LOAD\_STR \\
& 0x17 & LOAD\_BOOL \\
& 0x18 & LOAD\_DEC \\
\midrule
\multirow{7}{*}{Comparison} & 0xC0 & COMPARE \\
& 0xC3 & EQ \\
& 0xC4 & GE \\
& 0xC5 & GT \\
& 0xC6 & LE \\
& 0xC7 & LT \\
& 0xC8 & NE \\
\midrule
\multirow{4}{*}{Logical} & 0xD0 & AND \\
& 0xD1 & OR \\
& 0xD2 & NOT \\
& 0xD3 & XOR \\
\midrule
\multirow{6}{*}{Control Flow} & 0xD4 & LABEL \\
& 0xE4 & JUMP \\
& 0xE2 & CALL \\
& 0xE3 & RET \\
& 0xC1 & JUMP\_IF \\
& 0xC2 & CALL\_IF \\
\midrule
\multirow{2}{*}{I/O} & 0xE0 & PRINT \\
& 0xE1 & INPUT \\
\midrule
\multirow{4}{*}{Structural} & 0xF0 & ARC\_START \\
& 0xF1 & ARC\_END \\
& 0xF2 & ARC\_DELIM \\
& 0x00 & NOP \\
\bottomrule
\end{tabular}
\caption{Complete Opcode Reference}
\end{table}

\subsection{Keyword Reference}

\begin{table}[h]
\centering
\begin{tabular}{@{}ll@{}}
\toprule
Keyword & Description \\
\midrule
label & Define a label \\
let & Variable assignment \\
jump & Unconditional jump \\
goto & Alias for jump \\
call & Function call \\
return & Return from function \\
jumpif & Conditional jump \\
callif & Conditional call \\
exit & Terminate program \\
print & Output to stdout \\
input & Input from stdin \\
\bottomrule
\end{tabular}
\caption{Arcturus Keywords}
\end{table}

\subsection{Future Enhancements}

Potential areas for future development:

\begin{itemize}
    \item \textbf{Standard Library}: Collection of common functions (math, string manipulation, I/O)
    \item \textbf{Module System}: Namespace support for included files
    \item \textbf{Optimizations}: Constant folding, dead code elimination, peephole optimization
    \item \textbf{Type System}: Optional static typing for performance and safety
    \item \textbf{Arrays}: Native support for array data structures
    \item \textbf{FFI}: Foreign function interface for calling external code
    \item \textbf{Debugger}: Step-through debugger with breakpoints
    \item \textbf{REPL}: Interactive read-eval-print loop
\end{itemize}

\subsection{Design Decisions}

\subsubsection{Why Stack-Based?}

Stack-based VMs offer several advantages:
\begin{itemize}
    \item Simple instruction encoding
    \item Easy to understand and implement
    \item Compact bytecode
    \item Natural expression evaluation
\end{itemize}

\subsubsection{Why Global Variables?}

Global-only scope simplifies the implementation:
\begin{itemize}
    \item No need for scope tracking
    \item Simpler bytecode generation
    \item Easy to understand for beginners
    \item Explicit control over all state
\end{itemize}

Scoping can be added in future versions.

\subsubsection{Why Goto-Based Control Flow?}

Explicit labels and jumps provide:
\begin{itemize}
    \item Complete visibility into control flow
    \item Educational value (understanding low-level execution)
    \item Maximum flexibility
    \item Simple compilation from higher-level constructs
\end{itemize}

Higher-level constructs (if/while/for) can be added as syntactic sugar that compiles to labels and jumps.

\end{document}
